% =====================================================================
% apendice_a_checklist.tex -- Apendice A: Lista de verificacao
% consolidada do Capitulo 1, reunindo os criterios de aceitacao da
% NUREG-1537, Parte 2, ja distribuidos secao a secao no corpo do
% capitulo, com os criterios adicionais introduzidos pelo ISG-NUREG-1537
% Staff Guidance NPR-ISG-2011-002 (que complementa a NUREG-1537 para
% instalacoes de producao de radioisotopos e reatores homogeneos
% aquosos). Este e o primeiro dos tres apendices que documentam,
% respectivamente, os criterios de verificacao (A), a legislacao
% federal aplicavel (B) e a discussao sobre condicoes alem da base de
% projeto (C).
% =====================================================================
\nusecbreak{A}{Ap\^endice A}
\hypertarget{apA}{}%
\section*{Ap\^endice A -- Lista de Verifica\c{c}\~ao Consolidada}
\addcontentsline{toc}{section}{Ap\^endice A -- Lista de Verifica\c{c}\~ao Consolidada}
\label{sec:apendice-a}

Esta lista re\'une, em um \'unico local, todos os crit\'erios de
aceita\c{c}\~ao aplic\'aveis ao Cap\'itulo~1, tanto os da NUREG-1537,
Parte~2 \citep{nureg1537p2} -- j\'a apresentados, se\c{c}\~ao a se\c{c}\~ao,
nas tabelas de requisitos de aceita\c{c}\~ao do corpo deste cap\'itulo --
quanto os que o ISG-NUREG-1537 (\emph{Interim Staff Guidance Augmenting
NUREG-1537}, NPR-ISG-2011-002, \emph{``Final Interim Staff Guidance
Augmenting NUREG-1537 [\ldots] for Licensing Radioisotope Production
Facilities and Aqueous Homogeneous Reactors''})
acrescenta especificamente para instala\c{c}\~oes de produ\c{c}\~ao de
radiois\'otopos associadas a reatores n\~ao-de-pot\^encia -- caso do RMB,
conforme discutido no enquadramento regulat\'orio da Se\c{c}\~ao~1.1.

A coluna \textbf{Origem} identifica o documento-fonte de cada crit\'erio;
a coluna \textbf{ISG?} sinaliza rapidamente quais itens n\~ao t\^em
correspondente na NUREG-1537 original e foram introduzidos pelo ISG-NUREG-1537.
As colunas \textbf{Ref.\ no RPAS} e \textbf{Atendimento} ficam em aberto
nesta vers\~ao -- ser\~ao preenchidas, respectivamente, com a
localiza\c{c}\~ao exata no texto principal onde cada crit\'erio \'e
atendido e com um resumo objetivo de como ele \'e atendido, \`a medida
que o conte\'udo real da instala\c{c}\~ao for incorporado ao RPAS.

\begin{tabelaverificacao}
\multicolumn{6}{|l|}{\textbf{Se\c{c}\~ao 1.1 -- Introdu\c{c}\~ao}} \\
\hline
1.1-1 & O prop\'osito do RPAS deve estar claramente declarado. & NUREG-1537, Parte~2, \S1.1 & & Pendente & N\~ao \\
1.1-2 & O requerente deve estar identificado (nome, natureza jur\'idica). & NUREG-1537, Parte~2, \S1.1 & & Pendente & N\~ao \\
1.1-3 & A localiza\c{c}\~ao, o prop\'osito e o uso da instala\c{c}\~ao devem estar brevemente descritos. & NUREG-1537, Parte~2, \S1.1 & & Pendente & N\~ao \\
1.1-4 & As caracter\'isticas b\'asicas da instala\c{c}\~ao que afetam as considera\c{c}\~oes de licenciamento devem estar brevemente discutidas. & NUREG-1537, Parte~2, \S1.1 & & Pendente & N\~ao \\
1.1-5 & As caracter\'isticas de projeto ou de localiza\c{c}\~ao voltadas a endere\c{c}ar preocupa\c{c}\~oes b\'asicas de seguran\c{c}a devem estar delineadas. & NUREG-1537, Parte~2, \S1.1 & & Pendente & N\~ao \\
1.1-6 & Caracter\'isticas de seguran\c{c}a \'unicas do projeto, diferentes de instala\c{c}\~oes n\~ao-de-pot\^encia previamente licenciadas, devem estar destacadas. & NUREG-1537, Parte~2, \S1.1 & & Pendente & N\~ao \\
1.1-7 & \emph{Quando aplic\'avel:} as caracter\'isticas inerentes/passivas de seguran\c{c}a da instala\c{c}\~ao de produ\c{c}\~ao de radiois\'otopos associada devem tamb\'em estar descritas, n\~ao apenas as do reator. & ISG NPR-2011-002, \S1.1 & & Pendente & Sim \\
\hline
\multicolumn{6}{|l|}{\textbf{Se\c{c}\~ao 1.2 -- Resumo e Conclus\~oes sobre as Principais Considera\c{c}\~oes de Seguran\c{c}a}} \\
\hline
1.2-1 & Caracter\'isticas de projeto suficientes devem estar inclu\'idas para proteger a sa\'ude e a seguran\c{c}a do p\'ublico. & NUREG-1537, Parte~2, \S1.2 & & Pendente & N\~ao \\
1.2-2 & Nenhuma exposi\c{c}\~ao decorrente de opera\c{c}\~ao normal deve exceder os limites regulat\'orios aplic\'aveis, observado o princ\'ipio ALARA. & NUREG-1537, Parte~2, \S1.2 & & Pendente & N\~ao \\
1.2-3 & Os acidentes potenciais devem estar brevemente discutidos. & NUREG-1537, Parte~2, \S1.2 & & Pendente & N\~ao \\
1.2-4 & Todos os modos de opera\c{c}\~ao e eventos capazes de causar libera\c{c}\~oes radiol\'ogicas significativas e exposi\c{c}\~ao do p\'ublico devem estar discutidos. & NUREG-1537, Parte~2, \S1.2 & & Pendente & N\~ao \\
1.2-5 & \emph{Quando aplic\'avel:} a vis\~ao geral deve descrever a rela\c{c}\~ao entre as caracter\'isticas espec\'ificas da instala\c{c}\~ao de produ\c{c}\~ao e os principais processos ali conduzidos. & ISG NPR-2011-002, \S1.2 & & Pendente & Sim \\
1.2-6 & \emph{Quando aplic\'avel:} a descri\c{c}\~ao deve incluir a localiza\c{c}\~ao das edifica\c{c}\~oes dos principais componentes de processo, apoiada em desenhos de arranjo. & ISG NPR-2011-002, \S1.2 & & Pendente & Sim \\
1.2-7 & \emph{Quando aplic\'avel:} eventuais informa\c{c}\~oes propriet\'arias ou sens\'iveis relativas \`a instala\c{c}\~ao de produ\c{c}\~ao devem estar identificadas. & ISG NPR-2011-002, \S1.2 & & Pendente & Sim \\
1.2-8 & \emph{Quando aplic\'avel:} a vis\~ao geral de processo deve resumir os principais processos qu\'imicos/mec\^anicos envolvendo quantidades licenci\'aveis de material radioativo, com base, em parte, no resumo da an\'alise de seguran\c{c}a integrada (ISA). & ISG NPR-2011-002, \S1.2 & & Pendente & Sim \\
1.2-9 & \emph{Quando aplic\'avel:} o n\'ivel de detalhe da vis\~ao geral deve ser consistente com os resumos de ISA e com as an\'alises de acidentes do Cap\'itulo~13, ainda que menos detalhado. & ISG NPR-2011-002, \S1.2 & & Pendente & Sim \\
\hline
\multicolumn{6}{|l|}{\textbf{Se\c{c}\~ao 1.3 -- Descri\c{c}\~ao Geral}} \\
\hline
1.3-1 & Devem estar brevemente descritos: localiza\c{c}\~ao geogr\'afica, caracter\'isticas do s\'itio, crit\'erios de projeto, sistemas de seguran\c{c}a, instrumenta\c{c}\~ao/controle/el\'etrica, sistemas de refrigera\c{c}\~ao e auxiliares, gerenciamento de rejeitos e prote\c{c}\~ao radiol\'ogica, e instala\c{c}\~oes experimentais. & NUREG-1537, Parte~2, \S1.3 & & Pendente & N\~ao \\
1.3-2 & O arranjo geral das principais estruturas e equipamentos deve estar indicado por plantas e cortes (desenhos de eleva\c{c}\~ao). & NUREG-1537, Parte~2, \S1.3 & & Pendente & N\~ao \\
1.3-3 & As caracter\'isticas de seguran\c{c}a de maior interesse devem estar brevemente identificadas. & NUREG-1537, Parte~2, \S1.3 & & Pendente & N\~ao \\
1.3-4 & Caracter\'isticas incomuns do s\'itio, da cont\^enc\~ao, do projeto do reator ou de instala\c{c}\~oes experimentais \'unicas devem estar destacadas. & NUREG-1537, Parte~2, \S1.3 & & Pendente & N\~ao \\
\hline
\multicolumn{6}{|l|}{\textbf{Se\c{c}\~ao 1.4 -- Instala\c{c}\~oes e Equipamentos Compartilhados}} \\
\hline
1.4-1 & A instala\c{c}\~ao deve ser projetada para acomodar todos os usos ou falhas das instala\c{c}\~oes compartilhadas sem degrada\c{c}\~ao das caracter\'isticas de seguran\c{c}a do reator. & NUREG-1537, Parte~2, \S1.4 & & Pendente & N\~ao \\
1.4-2 & O projeto deve evitar condi\c{c}\~oes em que a contamina\c{c}\~ao possa se propagar para as instala\c{c}\~oes ou equipamentos compartilhados. & NUREG-1537, Parte~2, \S1.4 & & Pendente & N\~ao \\
1.4-3 & Onde necess\'ario, as barreiras devem estar brevemente descritas para garantir o atendimento aos dois crit\'erios anteriores. & NUREG-1537, Parte~2, \S1.4 & & Pendente & N\~ao \\
\hline
\multicolumn{6}{|l|}{\textbf{Se\c{c}\~ao 1.5 -- Compara\c{c}\~ao com Instala\c{c}\~oes Similares}} \\
\hline
1.5-1 & As compara\c{c}\~oes devem demonstrar que a instala\c{c}\~ao proposta n\~ao excede o envelope de seguran\c{c}a das instala\c{c}\~oes similares. & NUREG-1537, Parte~2, \S1.5 & & Pendente & N\~ao \\
1.5-2 & Deve haver garantia razo\'avel de que as exposi\c{c}\~oes radiol\'ogicas do p\'ublico n\~ao excedem os regulamentos e as diretrizes ALARA da instala\c{c}\~ao proposta. & NUREG-1537, Parte~2, \S1.5 & & Pendente & N\~ao \\
1.5-3 & \emph{Quando aplic\'avel:} a lista de instala\c{c}\~oes de refer\^encia deve incluir instala\c{c}\~oes ou opera\c{c}\~oes de produ\c{c}\~ao de radiois\'otopos compar\'aveis, al\'em dos reatores heterog\^eneos j\'a citados (os comparadores espec\'ificos de AHR listados pelo ISG-NUREG-1537 -- LOPO, HYPO, SUPO, TRACY, HRE -- n\~ao se aplicam ao RMB). & ISG NPR-2011-002, \S1.5 & & Pendente & Sim \\
\hline
\multicolumn{6}{|l|}{\textbf{Se\c{c}\~ao 1.6 -- Resumo das Opera\c{c}\~oes}} \\
\hline
1.6-1 & O requerente deve demonstrar a consist\^encia das opera\c{c}\~oes propostas com as premissas adotadas nos cap\'itulos posteriores do RPAS, incluindo os efeitos sobre a integridade do reator e as exposi\c{c}\~oes radiol\'ogicas potenciais. & NUREG-1537, Parte~2, \S1.6 & & Pendente & N\~ao \\
1.6-2 & O requerente deve demonstrar que a opera\c{c}\~ao proposta do reator foi considerada de forma conservadora no projeto e nas an\'alises de seguran\c{c}a. & NUREG-1537, Parte~2, \S1.6 & & Pendente & N\~ao \\
1.6-3 & (Renova\c{c}\~ao) As opera\c{c}\~oes propostas devem ser consistentes com as premissas dos cap\'itulos posteriores do RPAS. & NUREG-1537, Parte~2, \S1.6 & & Pendente & N\~ao \\
\hline
\multicolumn{6}{|l|}{\textbf{Se\c{c}\~ao 1.7 -- Conformidade com a Gest\~ao de Combust\'ivel Irradiado e Rejeitos}} \\
\hline
1.7-1 & (EUA) O requerente deve ter submetido um resumo do contrato com o DOE para a destina\c{c}\~ao de rejeitos de alta atividade e combust\'ivel irradiado. & NUREG-1537, Parte~2, \S1.7 & & N\~ao aplic\'avel fora dos EUA & N\~ao \\
1.7-2 & (EUA) O requerente deve ter indicado onde uma c\'opia da carta do contrato pode ser encontrada no RPAS. & NUREG-1537, Parte~2, \S1.7 & & N\~ao aplic\'avel fora dos EUA & N\~ao \\
1.7-3 & (Fora dos EUA -- adapta\c{c}\~ao) O requerente deve identificar o arcabou\c{c}o nacional aplic\'avel e demonstrar conformidade com ele. & S\'intese pr\'opria, a partir da NUREG-1537, Parte~2, \S1.7 & & Pendente & N\~ao \\
\hline
\multicolumn{6}{|l|}{\textbf{Se\c{c}\~ao 1.8 -- Modifica\c{c}\~oes e Hist\'orico da Instala\c{c}\~ao}} \\
\hline
1.8-1 & O requerente deve submeter um hist\'orico completo da instala\c{c}\~ao (quando aplic\'avel a renova\c{c}\~ao/emenda) ou das atividades pr\'evias relevantes (quando instala\c{c}\~ao nova). & NUREG-1537, Parte~2, \S1.8 & & Pendente & N\~ao \\
\end{tabelaverificacao}

\begin{notaregulatoria}{Sobre os itens marcados como oriundos do ISG-NUREG-1537}
Os itens sinalizados com ``Sim'' na coluna \textbf{ISG-NUREG-1537?} n\~ao existem
no texto original da NUREG-1537; foram introduzidos pelo NPR-ISG-2011-002
especificamente para instala\c{c}\~oes de produ\c{c}\~ao de radiois\'otopos
associadas a reatores n\~ao-de-pot\^encia -- cen\'ario aplic\'avel ao RMB
em raz\~ao da produ\c{c}\~ao de molibd\^enio-99, conforme discutido no
enquadramento regulat\'orio da Se\c{c}\~ao~1.1. Sua aplicabilidade efetiva
depende do escopo de cada fase de licenciamento: no escopo inicial deste
RPAS, limitado ao pr\'edio do reator, esses itens permanecem \emph{em
aberto} at\'e que a instala\c{c}\~ao de produ\c{c}\~ao de Mo-99 seja
formalmente incorporada.
\end{notaregulatoria}
